\chapter{Performance at a first glance: High-level, top-down analysis}
% \setcounter{section}{0}
\label{chapter:high-level}


There are various ways or strategies to obtain a first high-level overview and how to get started if we want to understand a code's performance better.
Every team, every project, every scientist have their own approach, and it often changes from project to project.
In this chapter, we propose a workflow running along few very simple principles:

\begin{enumerate}
  \item We study the performance characteristics for \emph{different machine part/aspects separately}. 
  \item We use a simple \emph{blackbox performance metric} per dimension to find out first if there are issues or not.
  \item We \emph{do not use any tools} (yet).
  \item We are not interested in explanations \emph{why} we see any performance flaws. We simple document \emph{what flaws} we observe.
\end{enumerate}

\noindent
Per dimension, we later break all performance flaws found down into subrubrics, and dig deeper and deeper into the code.
Along this way, we construct working hypotheses why the code performs poorly.
But this comes later.

